\section{An input order for Knuth's sixteen-pair graph}
\label{app:knuth16}

Answer~15(d) gives the graph but not an input order. Here is one under
which Algorithm~H follows the Tchoukaillon schedule and passes through
Knuth's matrices. We keep Knuth's logical names $g_k,b_k$. The edges are
$g_kb_k$ for $1\le k\le16$ and
\begin{equation}
\begin{gathered}
 b_1g_2,\ b_3g_4,\ b_5g_6,\ \ldots,\ b_{15}g_{16},\\
 b_2g_4,\ b_3g_5,\ b_8g_{10},\ b_9g_{11},\ b_{14}g_{16},\
 b_4g_6,\ b_5g_8,\ b_7g_9,\ b_6g_{10},\ b_3g_7,\ b_{10}g_{12},
\end{gathered}
\label{eq:knuth16edges}
\end{equation}
so $M=N=s=16$ and $E=35$.

\paragraph{The input.} Throughout this appendix $g_k$ and $b_k$ are Knuth's
logical names; the numbers that Algorithm~H sees as girl indices are a
separate thing, given in the column ``number'' below. Boys keep their
numbers. Girls receive the numbers $16,15,\ldots,1$ in the order
\[
  g_{16},g_{14},g_{12},g_{10},g_8,g_6,g_4,g_2,\ g_1,g_3,g_7,g_{13},g_9,g_{15},g_{11},g_5 .
\]
Each girl's list is in increasing order of boys, except that an even girl
$g_k$ puts $b_{k-1}$ first. In numerical form:
\begin{center}\small
\begin{tabular}{rrl@{\qquad}rrl}
\toprule
number&girl&list&number&girl&list\\
\midrule
1&$g_5$&$[3,5]$&9&$g_2$&$[1,2]$\\
2&$g_{11}$&$[9,11]$&10&$g_4$&$[3,2,4]$\\
3&$g_{15}$&$[15]$&11&$g_6$&$[5,4,6]$\\
4&$g_9$&$[7,9]$&12&$g_8$&$[7,5,8]$\\
5&$g_{13}$&$[13]$&13&$g_{10}$&$[9,6,8,10]$\\
6&$g_7$&$[3,7]$&14&$g_{12}$&$[11,10,12]$\\
7&$g_3$&$[3]$&15&$g_{14}$&$[13,14]$\\
8&$g_1$&$[1]$&16&$g_{16}$&$[15,14,16]$\\
\bottomrule
\end{tabular}
\end{center}

\paragraph{The run.} Step~\Hs1 scans the even girls first; each takes
$b_{k-1}$, so the matching is $b_1g_2$, $b_3g_4$, \dots, $b_{15}g_{16}$, as in
Knuth's first matrix, and~\Hs1 examines 20 entries. The rounds are:
\begin{center}\small
\begin{tabular}{crrrrrl}
\toprule
$L$&$f$&SAPs&$W_5$&$W_{13}$&$W_{11}$&SAPs found, as new couples $(g,b)$\\
\midrule
2&8&3&35&16&16&$(5,3),(4,2)$;\ $(11,9),(10,8)$;\ $(15,15),(16,14)$\\
3&5&2&34&14&15&$(9,7),(8,5),(6,4)$;\ $(13,13),(14,14),(16,16)$\\
4&3&1&28&10&12&$(7,3),(5,5),(8,8),(10,6)$\\
5&2&1&25&6&12&$(3,3),(7,7),(9,9),(11,11),(12,10)$\\
6&1&1&24&6&12&$(1,1),(2,2),(4,4),(6,6),(10,10),(12,12)$\\
\bottomrule
\end{tabular}
\end{center}
Here girls have their logical names. The free girls before the rounds are,
in \A{QUEUE} order, $(1,3,7,13,9,15,11,5)$, $(1,3,7,13,9)$, $(1,3,7)$,
$(1,3)$, $(1)$; every SAP ends at the last free girl, so step~\Hs{16} never
reorders the queue. After each round the basic paths with respect to the
matching $g_kb_k$ are exactly the rows of the next matrix in
Fig.~\ref{fig:knuth16rows}, and the run ends with that perfect matching.

\paragraph{All augmenting paths.} Enumerating every simple augmenting path
before each round gives the following counts, by level.
\begin{center}\small
\begin{tabular}{cl}
\toprule
before the round at&numbers of augmenting paths of each level\\
\midrule
$L=2$&22 of level~2, 4 of level~3\\
$L=3$&16 of level~3, 6 of level~5\\
$L=4$&7 of level~4, 12 of level~6, 8 of level~8\\
$L=5$&3 of level~5, 4 of level~7, 4 of level~9, 2 of level~11\\
$L=6$&1 of level~6\\
\bottomrule
\end{tabular}
\end{center}
So property $\Up$ fails in every round but the last, while property
$\Bp$ holds throughout. The paths $\pi_2$ and $\pi_3$ of~\eqref{eq:PQ} are two of
the 26 counted in the first line.
